% ============================================================================
%  第 4 章 风险建模 —— body/04-risk-model.tex
%  职责：把「风险的成本」拆成人能验证的子模块：概率 × 暴露 × 后果 × 集成
% ============================================================================

\section{时空异质性风险建模}
\label{sec:risk-model}

% ---------------------------------------------------------------------------
\subsection{时变失效概率}

沿用比例风险模型的构造，令基线危险率为 $\hazard_0(t)$，协变量向量
$\mathbf{z}(t)$ 包含风速、降雨与建筑峡谷因子，则
\begin{equation}
  \pcrashT{t}{\mathbf{x}, \mathbf{z}}
    = 1 - \exp\!\left[-\int_{0}^{t} \hazard_0(u)
      \exp\bigl(\boldsymbol{\beta}\transpose \mathbf{z}(u)\bigr) \, \dd{u}\right],
  \label{eq:hazard}
\end{equation}
其中 $\boldsymbol{\beta}$ 为回归系数向量，其物理意义为各协变量对
失效风险的弹性。

\todo{说明 $\hazard_0$ 的分布族选择（Weibull / 常数）与标定数据来源。}

% ---------------------------------------------------------------------------
\subsection{暴露度演化}

地表 $h$ 处单元 $m$ 在时刻 $t$ 的暴露密度由静态基底与活动强度合成：
\begin{equation}
  \exposure(m, t) = \exposure_{\mathrm{static}}(m)
    + \exposure_{\mathrm{tidal}}(m) \, \phi\bigl(t; \pi(m)\bigr),
  \label{eq:exposure}
\end{equation}
其中 $\pi(m)$ 标识土地利用类型，$\phi(\cdot)$ 为对应的昼夜激活曲线。

% ---------------------------------------------------------------------------
\subsection{后果代价：致死、财产与噪声}

三个维度的后果代价分别为 $\casualty$、$\propertyloss$ 与 $\noisecost$。
以致死代价为例：
\begin{equation}
  \casualty(m, t) = \exposure(m, t) \, A_{\mathrm{exp}}
    \, \pcrashT{t}{\mathbf{x}_m, \mathbf{z}}.
  \label{eq:casualty}
\end{equation}

% ---------------------------------------------------------------------------
\subsection{综合集成与性质}

\begin{definition}[综合风险代价]
给定边 $(u \to v)$ 与其通过时刻 $t$，定义
\begin{equation}
  \riskcost\bigl((u, t) \to (v, t')\bigr)
    = w_1 \casualty + w_2 \propertyloss + w_3 \noisecost,
  \quad w_i \ge 0,\; \sum_i w_i = 1.
  \label{eq:risk-cost}
\end{equation}
\end{definition}

\begin{theorem}[单调性]
\label{thm:monotone}
若 $\hazard_0(t)$ 与 $\boldsymbol{\beta}\transpose \mathbf{z}(t)$ 均非负有界，
则由 \refeqn{eq:risk-cost} 定义的边代价满足
$\riskcost(t_2) \ge \riskcost(t_1)$ 对 $t_2 > t_1$ 未必成立，
但其积分形式对时间可测且满足 Bellman 最优性方程的前提条件。
\end{theorem}

\begin{proof}
证明见 \refapp{app:supplement}。
\todo{补完整证明：核心是证明可加性 + FIFO + 非负性三条。}
\end{proof}
